\section{Persistance de voie, multisections d'espèce et conjugaison de Möbius}
\label{sec:persistencia-ruta-mobius-familias}

Cette section précise une distinction qui, formulée de manière trop
abrégée, inverse la logique de la construction. La voie n'est pas une donnée qu'une
particule choisit une fois sa masse connue. Ce n'est pas non plus un dessin auxiliaire
sans contenu causal. C'est l'histoire observable et orientée d'une extension
survivante. L'extension conserve, en outre, la profondeur compatible, le
registre, les retenues et la mémoire qui ne tiennent pas dans une coupe finie. En ce
sens précis,
\[
 \boxed{\text{une particule HMT est la persistance compatible d'une route
 enrichie}.}
\]
La phrase et la définition par extensions expriment le même objet à deux
niveaux : la voie est son histoire visible ; l'extension est la tour complète de
cette histoire avec tous ses prolongements compatibles.

\subsection{Persistance comme section d'un système inverse}

Soit \((\mathcal S_m,\rho_{m+1,m})_{m\geq0}\) le système d'états qui
survivent jusqu'à la profondeur \(m\), et soit
\[
 x=(x_m)_{m\geq0}\in\Omega_{\rm HMT}^{\rm surv}
 =\varprojlim_m\mathcal S_m.
\]
Les ombres de voie satisfont le carré de compatibilité
\[
 \operatorname{tr}_{m+1,m}\!\left(\operatorname{sh}_{m+1}(x_{m+1})\right)
 =\operatorname{sh}_{m}\!\left(\rho_{m+1,m}(x_{m+1})\right)
 =\operatorname{sh}_{m}(x_m).
\]
Par conséquent
\[
 \gamma_x=\bigl(\gamma_x^{[m]}\bigr)_{m\geq0},
 \qquad \gamma_x^{[m]}=\operatorname{sh}_m(x_m),
\]
est une histoire compatible, non une séquence ajustée rétrospectivement.

Sur chaque coupe \(\gamma_x^{[m]}\), on dispose du fibré discret
\(\mathcal B_{\gamma_x^{[m]}}\). Une section persistante est une famille
\[
 s=(s_m)_{m\geq0},
 \qquad
 s_m\in\Gamma\!\left(\mathcal B_{\gamma_x^{[m]}}\right),
 \qquad
 \operatorname{res}_{m+1,m}(s_{m+1})=s_m.
\]
Cette équation est le contenu mathématique de la persistance : chaque lecture
finie se prolonge sans contredire les lectures précédentes. Le tuple de
particule du \cref{chap:particula} ajoute l'occupation,
la représentation et la propagation, mais n'altère pas cette provenance.

\subsection{Espèce interne comme quotient orbital de la multisection}

Dans l'atlas nonadique, chaque cellule orientée \(c=(r,c')\in\mathcal C_{81}\)
porte une famille de voie
\[
 \mathcal R(c)=
 (n_A,s_C,k_{\Delta_4},\nu_{120},\nu_{270},\tau_{12}).
\]
L'image de \(\mathcal R\) contient cinquante-six valeurs. Si l'on fixe la
carte et son ordre, le rang \(\kappa\in\{0,1,2,3,4\}\) au sein de chaque fibre
produit le relèvement injectif
\[
 c\longmapsto \bigl(\mathcal R(c),\kappa(c)\bigr).
\]
L'alphabet à cinq états est minimal, car il existe deux fibres de
cardinal cinq. Cette sélection est interne, prospective et aveugle aux masses et aux
noms de particules.

Le recensement complet qui soutient cette affirmation est
\[
 81=45_{\rm leptones}+36_{\rm quarks}
   =25_{\ell^-}+20_{\nu}+20_{d\text{-type}}+16_{u\text{-type}}.
\]
Les cinquante-six fibres de voie possèdent l'histogramme
\[
 41\times1+10\times2+2\times3+1\times4+2\times5=81.
\]
Par conséquent, \(56\) compte les signatures de voie, tandis que \(81\) compte les cellules
orientées ; la mémoire \(\kappa\) distingue les cellules qu'une même signature ne
sépare pas.

Pour classer les espèces internes, on n'inverse pas cette application. On prend le
quotient
\[
 q:\mathcal C_{81}\longrightarrow\Sigma_{13},
 \qquad
 q(c)=\bigl(\operatorname{familia}(c),G(c)\bigr),
\]
dont l'image comprend treize types. La sortie associée à un type \(y\) est la
multisection finie de l'atlas
\[
 \boxed{
 \mathfrak S_y^{\rm atlas}=
 \{(\mathcal R(c),\kappa(c)):q(c)=y\}.}
\]
Ses cardinaux sont
\[
 1,8,16;\qquad 2,4,6,8;\qquad 2,4,6,8;\qquad4,12,
\]
pour les leptons chargés, les neutrinos et les deux types de quark,
respectivement. L'espèce est le type commun de cette fibre ; une particule
concrète est une section persistante réalisée en son sein. Exiger que le
nom d'une espèce choisisse à lui seul une cellule détruit précisément les
degrés d'orientation et de conjugaison que la fibre conserve.

Le relèvement aux extensions est un objet typé différent. Si
\(\pi_{81}:\Omega_{\rm HMT}^{\rm surv}\to\mathcal C_{81}\) désigne le lecteur
cellulaire de la voie enrichie déclarée, alors
\[
 \widetilde{\mathfrak S}_y
 =\{x\in\Omega_{\rm HMT}^{\rm surv}:q(\pi_{81}x)=y\}
\]
est la multisection enrichie correspondante. Le certificat fini
démontre le recensement et l'obstruction dans
\(\mathfrak S_y^{\rm atlas}\) ; son relèvement est relatif au lecteur
\(\pi_{81}\), non une identification automatique entre une cellule et une
histoire profonde.

\begin{theorem}[Obstruction au sélecteur ponctuel et exception électronique]
Soit \(\rho_{\rm c}(r,c')=(10-r,10-c')\) l'inversion centrale. Les treize fibres de
\(q\) sont \(\rho_{\rm c}\)-stables. La fibre de lepton chargé de niveau zéro possède
l'unique point fixe, \((5,5)\). Aucune des douze autres fibres n'admet
un sélecteur ponctuel \(\rho_{\rm c}\)-équivariant si \(\rho_{\rm c}\) agit trivialement
sur le type grossier.
\end{theorem}
\begin{proof}
L'équation \(\rho_{\rm c}(r,c')=(r,c')\) équivaut à \(r=c'=5\). Le recensement de l'atlas
situe cette cellule dans la fibre \((\ell^-,G0)\), qui a pour cardinal un. Soit
maintenant \(y\neq(\ell^-,G0)\) et supposons qu'il existe une section ponctuelle
équivariante \(a:\Sigma_{13}\to\mathcal C_{81}\). Comme \(\rho_{\rm c} y=y\),
l'équivariance exigerait
\[
 a(y)=a(\rho_{\rm c} y)=\rho_{\rm c} a(y).
\]
Alors \(a(y)\) devrait être un point fixe de \(\rho_{\rm c}\), mais l'unique
point fixe appartient à la fibre électronique. C'est impossible. Par conséquent,
dans les douze autres fibres, la sortie naturelle est une \(\rho_{\rm c}\)-orbite ou une
multisection ; choisir un représentant requiert une orientation ou une
représentation supplémentaire.
\end{proof}

L'électron est exceptionnel en trois sens compatibles : \((5,5)\) est le
centre géométrique unique, l'unique cellule de niveau zéro et le centre par rapport
auquel se forme le réseau brut de différences
\[
 \Lambda_{\rm raw}^{\Delta}
 =\bigl\langle A_{\rm raw}(c)-A_{\rm raw}(5,5):
 c\in\mathcal C_{81}\bigr\rangle_{\mathbb Z}.
\]
La signature brute de cette cellule n'est pas nulle :
\[
 A_{\rm raw}(5,5)=(24,0,12,0,-12),
 \qquad \tau(5,5)=\texttt{+++---+++---}.
\]
L'origine \(v_e=(0,0,0,0,0)\) du réseau d'action est une convention
relative à l'électron et non la signature brute de la cellule centrale.

L'invariant algébrique électronique est évalué au moyen de
\[
 \mathfrak e_0=
 \frac{\sqrt3}{4}
 \left(e^{\varphi/\pi^2}-22\alpha_{\rm HMT}^{3}\right)
 (1+15\Delta_4),
 \qquad
 \Delta_4=
 \frac{\pi-e\log\pi}{270}\left(1+\frac{\pi}{729}\right),
\]
\[
 \mathfrak e_0
 =0.510998922488066072509870877191349\ldots.
\]
L'expression est adimensionnelle ; la carte
\(\varsigma_{u_E}:y\mapsto y\,u_E\) introduit ensuite une unité d'énergie.
Le centre \((5,5)\) et son unicité sont des résultats finis exacts. La sélection
des coefficients \(22\) et \(15=3\cdot5\) se déduit, avec la structure de
départ explicite, au moyen du complément
des cinq régions de \(\pi\) dans \(K_e=\mathbb Z/9\mathbb Z\times\mathbb F_3\) :
\(-22=-|K_e\setminus\iota(R_\pi)|\) et
\(+15=|R_\pi\times\mathbb F_3|\). La signature brute centrale conserve un type
distinct et agit comme contrôle indépendant de la sélection basale.

La cellule \((5,5)\) ne doit pas être confondue avec la signature énergétique
\(555555\) : cette dernière appartient au représentant transversal d'une des
cinq régions de \(\pi\). Ce sont des objets distincts qui partagent le chiffre
central.

Dans le secteur des quarks, le caractère résiduel
\[
 \operatorname{col}(r,c')=r-c'\pmod3
\]
raffine la multisection en trois classes de douze cellules et satisfait
\(\operatorname{col}(\rho_{\rm c} c)=-\operatorname{col}(c)\). Cela construit le
tritique interne de couleur. La représentation physique de \(SU(3)\) est un
lecteur ultérieur ; elle n'est pas nécessaire pour démontrer la répartition finie.

\subsection{Stratification d'extension, de voie, de registre et de signature}

Les applications \(q\), \(\mathcal R\), la projection cellulaire et le lecteur de
réalisation distinguent quatre classifications liées et typées :
\begin{enumerate}[label=(\roman*)]
 \item les treize fibres internes famille--niveau de \(q\) ;
 \item les cinquante-six familles de voie de \(\mathcal R\) ;
 \item les quatre-vingt-une positions et leurs trois cent vingt-quatre
       orientations ;
 \item les réalisations physiques individualisées par représentation et codomaine.
\end{enumerate}
Les quatre couches sont recensées avant la confrontation métrologique. Pour chaque classe,
la comparaison requiert le tuple
\[
(\text{position},\text{fibre},\text{routes},\text{lecteurs},
\text{registre},\text{signature},\text{caractère},\text{tours},
\text{carte dimensionnelle}).
\]
Une reconstruction sur un sous-ensemble démontre seulement ce sous-cas et ne
remplace aucune des composantes du tuple.
Étant donnés \(A\), \(C^*\) et les canaux \(q_\pm\), la tour globale appartient au
semi-groupe infini
\[
 s_n=q_-^n-q_+^n,
 \qquad
 n\in\langle90,120\rangle
 =\{90,120\}\cup\{30r:r\geq6\}.
\]
Les hauteurs \(90,120,180,210,240,270,360,420\) sont des représentants finis du
semi-groupe ; la récurrence précédente prolonge la tour complète.
et le caractère adimensionnel d'une signature est
\[
 \chi_{\rm mass}(v)=\exp\!\left(
 n_AA_{\rm rad}+\frac{s_C}{6}C^*_{\rm rad}
 +k\Delta_4+\nu_{120}s_{120}+\nu_{270}(135\Delta_4^2)
 \right).
\]
Les moments sont globaux : les espèces les pondèrent au moyen de leur voie, de leur
registre et de leur opérateur de fibre. Les lecteurs \(K_D\) et \(K_\Omega\)
déterminent le raffinement bidirectionnel avant la confrontation, tandis que la
table intégrale des espèces publie séparément la généalogie interne, la
carte dimensionnelle et la comparaison externe.

\subsection{Conjugaison, anti-section et ruban de Möbius}

Sur un mot dodécaphasique, on définit l'involution exacte
\[
 C_M(\tau)_m=-\tau_{11-m},
 \qquad C_M^2=I.
\]
Les fonctionnelles linéaires à poids symétriques sont impaires sous \(C_M\),
tandis que les corrélations quadratiques compatibles avec le renversement sont
paires. Ainsi, l'orientation et la charge linéaire vivent dans le secteur impair,
tandis que la mémoire quadratique vit dans le secteur pair. Cela ne démontre pas que
le caractère exponentiel complet de masse soit pair : si
\(\chi(\tau)=e^{\ell(\tau)}\) et que \(\ell\) est impair, alors
\(\chi(C_M\tau)=\chi(\tau)^{-1}\).

L'involution précédente agit exactement sur les mots. Pour parler
d'anti-section, on fixe en outre un relèvement \(\widetilde C_M\) au fibré des
sections qui préserve l'admissibilité, le registre et les troncatures :
\[
 \operatorname{res}_{m+1,m}\widetilde C_M
 =\widetilde C_M\operatorname{res}_{m+1,m},
 \qquad
 \operatorname{sh}(\widetilde C_Mx)=C_M\operatorname{sh}(x).
\]
Ce n'est que sous ce relèvement que l'anti-section de \(s\) est
\(\widetilde C_Ms\) sur la voie inversée.

Le mot ``Möbius'' possède ici une formulation relative précise. Soit
\(\ell_{\rm ret}\) une boucle de retour de la voie enrichie qui engendre
\[
 \langle\ell_{\rm ret}\rangle\simeq\pi_1(S^1)\simeq\mathbb Z.
\]
Ce groupe de retour n'est pas la phase finie \(C_9\). Supposons que son
relèvement au
registre possède un caractère de signe
\[
 \varepsilon_x:\langle\ell_{\rm ret}\rangle\longrightarrow\{\pm1\},
 \qquad \varepsilon_x(\ell_{\rm ret})=-1.
\]
Le fibré en intervalles associé est
\[
 \mathcal M_x=
 ([0,1]\times[-1,1])/(0,u)\sim(1,-u).
\]
Sa première classe de Stiefel--Whitney est non triviale ; par conséquent
\(\mathcal M_x\) est un ruban de Möbius. Un tour échange les deux
feuillets locaux et deux tours restaurent l'orientation. Cela n'impose pas que
le registre complet revienne au même état : retour d'orientation et retour
de mémoire demeurent des types différents.

Le théorème est relatif au relèvement, à la boucle de retour et au caractère
\(\varepsilon_x\), tous déclarés. Il ne naît pas de l'observation de deux signes.
La lecture des feuillets requiert, en outre, de distinguer la graduation scindée \(R\)
de la conjugaison \(C\) et d'imposer
\[
 CRC^{-1}=-R,
 \qquad CT(C^*)C^{-1}=T(-C^*).
\]
La lecture particule--antiparticule exige finalement que la représentation
physique inverse les porteurs impairs, préserve les porteurs quadratiques
pairs et entrelace l'évolution et le caractère complet de la manière déclarée
par la représentation. Sous ces prémisses, particule et antiparticule sont
les deux orientations conjuguées d'une même bande de persistance.

\begin{theorem}[Décomposition paire--impaire sur un revêtement double]
\label{pf84:A10-md-mobius}
Soit \(p:\widetilde X\to X\) un revêtement principal double d'involution
\(\iota\). Pour toute \(f\in C^0(\widetilde X,\mathbb R)\), les projecteurs
\[
 f^+=\frac12(f+\iota^*f),\qquad
 f^-=\frac12(f-\iota^*f)
\]
donnent la décomposition \(f=f^++f^-\), avec
\[
 f^+\circ\iota=f^+,\qquad f^-\circ\iota=-f^-.
\]
La partie paire descend à une fonction sur \(X\) ; la partie impaire est une
section du fibré en droites de signe associé. Pour le revêtement d'orientation
du ruban de Möbius, \(w_1\ne0\) fait obstruction à une section globale continue
et non nulle de ce fibré en droites.
\end{theorem}

\begin{proof}
Les identités des deux projecteurs prouvent la somme et les parités. Une
fonction invariante est constante sur les fibres de \(p\), donc
descend ; une fonction anti-invariante satisfait la loi de transition de la
représentation de signe. Dans le ruban de Möbius, le fibré en droites associé est non
trivial et sa première classe de Stiefel--Whitney est non nulle, de sorte que toute
section continue globale possède un zéro. La conclusion exclut une
section globale continue \emph{sans zéros} ; elle ne nie pas l'existence de
sections avec des zéros.
\end{proof}

Ce lemme classique type le secteur pair--impair du revêtement ; il n'identifie pas
à lui seul le renversement de mot, l'holonomie de retour, la
graduation scindée, la conjugaison des feuillets ni la parité d'une signature de
masse. Ces identifications requièrent leurs propres opérateurs d'entrelacement.

\subsection{TRIT algébrique et lecteur temporel}

La réalisation quadratique déclarée associe au signe visible du TRIT un type
algébrique. Pour
\(\bar\tau\in\{+1,0,-1\}\), soit
\[
 \mathcal A_{\bar\tau}=
 \mathbb R[J]/(J^2+\bar\tau).
\]
Alors
\[
 \begin{array}{c|c|c}
 \bar\tau&J^2&\text{régime}\ \\ \hline
 +1&-1&\text{elliptique/complexe}\ \\
 0&0&\text{parabolique/nilpotent}\ \\
 -1&+1&\text{hyperbolique/scindé}.
 \end{array}
\]
Les exponentielles correspondantes sont
\[
 e^{tJ}=\cos t+J\sin t,
 \qquad e^{tJ}=1+tJ,
 \qquad e^{tJ}=\cosh t+J\sinh t,
\]
et dans le secteur scindé apparaissent les projecteurs
\(P_\pm=(1\pm J)/2\). Ces identités se vérifient exactement dans la
réalisation déclarée. La
courbure ne provient pas du signe isolé : elle apparaît lorsque APP apporte des liens et
des deux-cellules et que l'on fixe une connexion. Dans la polarisation mixte somme--produit
publiée, sur les \(81\) plaquettes, la courbure orientée satisfait
\[
 F(S,S)=F(P,P)=0,
 \qquad F(S,P)=+1,
 \qquad F(P,S)=-1\pmod9.
\]
L'identité de Stokes finie correspondante contrôle que le bord de chaque
rectangle reproduit la somme de ses plaquettes.

En ordonnant les transitions, on obtient le lecteur temporel tritique :
\[
 \begin{aligned}
 +1&\longmapsto\text{retour, mémoire et passé},\\
  0&\longmapsto\text{seuil torsionnel et présent},\\
 -1&\longmapsto\text{ouverture bidirectionnelle et futur}.
 \end{aligned}
\]
Son précurseur exact est la différence entre retour de phase et retour
d'état :
\[
 g(k+9)=g(k),
 \qquad B_{k+9}\neq B_k\quad\text{en général}.
\]
Dans la réalisation géométrique correspondante, on parle de secteur
sphérique--volumétrique, de seuil plan avec registre torsionnel et de secteur
hyperbolique--surfacique. Cette correspondance définit un lecteur de mécanique
dimensionnelle soumis à l'application entre le TRIT et la géométrie ; elle ne redéfinit pas le TRIT
algébrique et n'identifie pas à elle seule
un signe à une courbure de l'espace-temps. L'expression ``cristal
temporel'' nomme cette architecture dans la terminologie HMT ; une réalisation
dynamique de cristal temporel exige en outre un hamiltonien ou un protocole
périodique d'excitation,
une réponse sous-harmonique robuste et une stabilité.

\subsection{Support exact de la lecture aller--retour}

L'échange bidirectionnel ne repose pas seulement sur une image verbale. Les
coordonnées angulaires entières
\[
 W_+=6n_A+s_C,
 \qquad W_-=6n_A-s_C
\]
s'obtiennent au moyen de
\[
 \binom{W_+}{W_-}
 =\begin{pmatrix}6&1\\6&-1\end{pmatrix}
 \binom{n_A}{s_C},
 \qquad
 \operatorname{SNF}\!\begin{pmatrix}6&1\\6&-1\end{pmatrix}
 =\operatorname{diag}(1,12).
\]
Le réseau est donc d'indice douze. L'involution
\(s_C\mapsto-s_C\) échange \(W_+\leftrightarrow W_-\) et, de manière
équivalente, \(q_+\leftrightarrow q_-\). Tel est le substrat algébrique exact
de l'aller et du retour. Son interprétation comme propagation avancée et retardée
requiert l'opérateur d'évolution de la représentation physique.

La projection paire employée dans la lecture biangulaire satisfait, pour toute
fonction pour laquelle les expressions sont définies,
\[
 \frac{f(\phi)+f(-\phi)}2=f_{\rm par}(\phi),
 \qquad
 \frac{(\pi-\phi)^2+(\pi+\phi)^2}{2}=\pi^2+\phi^2.
\]
Ce sont des identités exactes de symétrisation. La lecture en termes d'une
théorie physique de l'absorbeur relève du foncteur de propagation, non de
l'identité algébrique isolée.

\subsection{Périodicités spinorielles non équivalentes de \texorpdfstring{\(2\pi\)}{2 pi} et \texorpdfstring{\(4\pi\)}{4 pi}}

Dans le relèvement spinoriel classique,
\[
 U(\theta)=\exp\!\left(-\frac{i\theta}{2}\sigma_z\right),
 \qquad U(2\pi)=-I,
 \qquad U(4\pi)=I.
\]
L'égalité est exacte dans le revêtement double \(SU(2)\to SO(3)\). Son application
à une section HMT requiert de spécifier la représentation spinorielle du
transport. Une fois fixée, elle explique pourquoi un tour change le signe de la
section et deux tours le restaurent. Si l'on adopte les relations CAR dans
l'algèbre d'occupation, l'exclusion s'en déduit ; ni CAR ni Pauli ne
s'infèrent du signe de \(2\pi\) sans cette structure supplémentaire.

L'interprétation HMT ``feuillet surfacique \(2\pi\), feuillet volumétrique \(4\pi\)'' est
un autre objet. La section~\ref{sec:bisagra-hojas-pi-phi2} formalise son noyau
topologique : un tour de la base ferme le bord projeté en \(2\pi\) et
se termine sur le feuillet conjugué, tandis que le carré de la boucle ferme le revêtement
orienté en \(4\pi\). Le théorème est relatif au relèvement, à l'holonomie et au
lecteur géométrique déclarés. Son identification supplémentaire à une
représentation spinorielle requiert toujours le foncteur physique et ne se déduit pas
du partage des mêmes facteurs.

\subsection{Application entre multisections internes et représentations physiques}

L'objet particule, la projection
extension--voie, le sélecteur interne des 81 cellules, le centre électronique,
les treize multisections finies de famille, leur relèvement relatif au lecteur
cellulaire \(\pi_{81}\), le tritique de couleur, le registre, la signature, le caractère et
les tours forment la composition
\[
\text{extension}\longrightarrow\text{route}\longrightarrow
\text{multisection}\longrightarrow\text{signature}\longrightarrow
\text{caractère}\longrightarrow\text{tours}.
\]

L'obligation restante consiste à construire un foncteur de
réalisation qui apparie certaines étiquettes physiques individualisées et
composées aux multisections enrichies pertinentes et reproduise les
signatures raffinées déclarées sans consulter les masses ni les signatures cibles. Pour les
états composés, le registre est composé avant d'être projeté sur cinq
coordonnées. Le domaine du foncteur conserve l'électron central, l'atlas des
familles et la loi interne de caractère.


\paragraph{Typage probatoire de la persistance, de l'atlas et des familles.}
La persistance de voie, les anti-sections, la lecture bidirectionnelle et le
lecteur temporel tritique fixent l'architecture de départ. Le système
inverse, le sélecteur interne et la séparation des types réalisent cette
architecture au moyen de morphismes déclarés. Les recensements de l'atlas et les
identités de \(C_M\) sur les mots sont satisfaits dans tout le système fini ;
leur relèvement aux sections et le ruban de Möbius sont démontrés avec une
structure de départ explicite. Les coefficients électroniques
\((-22,+15)\) se déduisent avec la même structure explicite. Si un
relèvement ouvre au préalable une différence métrologique ou une signature cible,
sa valeur est une reconstruction calibrée.
